% GDZ PPN236018647,native559lower–561;HKUST560–562.
\NumberedItem{87}
\indent Series
\begingroup
\makeatletter\tagsleft@true\let\veqno\@@leqno\makeatother
\begin{equation}
1+\frac14+\left(\frac{1\,.\,1}{2\,.\,4}\right)^2+\left(\frac{1\,.\,1\,.\,3}{2\,.\,4\,.\,6}\right)^2+\text{etc.}=\frac4\pi\tag*{[1]}
\end{equation}
\par\endgroup
\noindent simul cum theoria generali serierum involventium sinus et cosinus angulorum arithmetice crescentium.
\par\noindent\hfill[1798] Iun.

\medskip\begingroup\small
\renewcommand{\thefootnote}{\ensuremath{\dagger})}
\indent Welche Teile der allgemeinen Lehre von den trigonometrischen Reihen \textsc{Gauss} im Auge hatte, als er die vorstehende Notiz schrieb, ergibt sich aus den Anfangsworten des oben S. 433 abgedruckten Briefes an \textsc{Schumacher} vom 5. Februar 1850. Danach handelte es sich um die »Methode, den Grad der Konvergenz« einer solchen Reihe zu bestimmen, d. h. also um die Untersuchung des asymptotischen Verhaltens des allgemeinen Gliedes für große Werte des Stellenzeigers, und zwar insbesondere bei der Entwicklung der Mittelpunktsgleichung nach den Sinus der Vielfachen der mittleren Anomalie. Wir wissen ja auch z. B. aus den Bemerkungen zu \textsc{Ulugh-Beighs} Tafeln\footnote{Diese Bemerkungen hatte \textsc{Gauss} im Frühjahr 1799 brieflich an v. \textsc{Zach} mitgeteilt, vergl. die genauen Angaben oben S. 444 in den Bemerkungen zum \textit{Briefwechsel} der \textit{Reihenlehre}.}, daß diese Untersuchungen von \textsc{Gauss} bis in die letzten Jahre des achtzehnten Jahrhunderts zurückreichen. Das Studium der Schriften von \textsc{Laplace}, mit dem
\par\endgroup

\clearpage\SourceLeaf{560}\ArithmeticHead{531}{R}{1798 MAI.—1798 IUN.}
\begingroup\small
\makeatletter\tagsleft@true\let\veqno\@@leqno\makeatother
\noindent \textsc{Gauss} nach dem Zeugnis der Tagebuchnotizen Nr. 86 und Nr. 88 um die hier in Rede stehende Zeit beschäftigt war, ist auch für diese Untersuchungen von Bedeutung gewesen.

Mit diesen geschichtlichen Feststellungen läßt sich auch das Auftreten der Reihe [1] in Einklang bringen. Diese Reihe für \(\frac4\pi\) hat \textsc{Gauss} auch in seinen Abdruck des zweiten Bandes von \textsc{Schulzes} Tafelwerk eingeschrieben und zwar mit dem Zusatz »\textsc{Euler}, H. d. l'A. 1778, p. 609«. In der Tat findet man in der Histoire de l'Académie des Sciences für das Jahr 1778 (Paris 1781), beginnend auf S. 603 der Mémoires: \textit{Extraits de différentes lettres de M. Euler à M. le Marquis de Condorcet}, und in dem dritten Briefe (vom 12./23. September 1776), der auf S. 606 beginnt, wird S. 609 (L. \textsc{Euleri} Opera omnia, ser. I, vol. 18, S. 77) die Reihe [1] durch Betrachtung von Integralen der Form
\[
\int_0^1 dx\left(\log\frac1x\right)^\lambda=\int_0^\infty u^\lambda e^{-u}\,du,
\]
\noindent also von \(\Pi\)-Funktionen (in der Bezeichnung der \textit{Disquisitiones circa seriem}, 1812, Werke III, S. 151), gewonnen. Da bei dieser Herleitung trigonometrische Reihen nicht ins Spiel kommen, so muß \textsc{Gauss}, als er die Tagebuchnotiz schrieb, ein von dem \textsc{Euler}schen abweichendes Verfahren im Sinn gehabt haben. Dazu kommt, daß in \textsc{Schulzes} \textit{Tafeln} neben der Reihe [1] noch die Reihe
\begin{equation}
1+\frac{1\,.\,1}{2\,.\,4}+\frac{1\,.\,1\,.\,3\,.\,3}{2\,.\,4\,.\,4\,.\,6}+\frac{1\,.\,1\,.\,3\,.\,3\,.\,5\,.\,5}{2\,.\,4\,.\,4\,.\,6\,.\,6\,.\,8}+\cdots\text{ in inf.}=\frac4\pi\tag*{[2]}
\end{equation}
\noindent eingeschrieben ist; \textsc{Gauss} dürfte also auf \textls[120]{beide} Reihen durch ein einheitliches Verfahren gekommen sein, das auf einer Entwicklung in eine trigonometrische Reihe beruhte. Im Nachlaß findet sich keine Aufzeichnung, die hierüber Aufklärung gibt, man ist also auf Vermutungen angewiesen. Nun hat \textsc{James Ivory} in demselben Jahre, aus dem auch unsere Tagebuchnotiz stammt, eine Abhandlung veröffentlicht\footnote{J. \textsc{Ivory}, \textit{A new series for the rectification of the Ellipsis; together with some observations on the Evolution of the Formula} \((a^2+b^2-2ab\cos\Phi)^n\), Transactions of the Royal Society of Edinburgh \textbf{4}, 1798, Part II, S. 177—190. Auf diese wenig bekannte Abhandlung verweist H. \textsc{Burkhardt}, \textit{Encyklopädie der mathem. Wissenschaften} II A 12, S. 883.}, in der jene \textsc{Euler}sche Reihe [1] nach einer Methode hergeleitet wird, die sich in einem \textsc{Gauss} durchaus vertrauten Gedankenkreise bewegt, und die mit den nötigen Abänderungen auch zur Ableitung der Reihe [2] benutzt werden kann. Sie gründet sich auf die von \textsc{Lagrange}\begingroup\renewcommand{\thefootnote}{**)}\footnote{J. L. \textsc{Lagrange}, \textit{Solution de différents problèmes de calcul intégral}, Miscell. Taurin. \textbf{3}, 1762—1765, Oeuvres I, Paris 1867, S. 469, insbesondere S. 620 ff.; \textit{Recherches sur les inégalités des Satellites de Jupiter}, Prix de l'Académie des Sciences \textbf{9}, Paris 1766, Oeuvres VI, Paris 1873, S. 63, insbesondere S. 88.}\endgroup gegebene trigonometrische Entwicklung
\begin{equation}
\begin{aligned}
(a^2+b^2-2ab\cos\theta)^n&=(a-be^{\theta\sqrt{-1}})^n(a-be^{-\theta\sqrt{-1}})^n\\
&=P+2Q\cos\theta+2R\cos2\theta+\cdots\text{ in inf.},
\end{aligned}\tag*{(3)}
\end{equation}
\noindent deren Koeffizienten \(P,Q,R,\ldots\) in der Form von Potenzreihen der Größe \(\frac ba\) dargestellt werden. Diese Entwicklung war namentlich ihrer astronomischen Anwendungen wegen sehr bekannt und berühmt; \textsc{Gauss} wird sie 1798 ebensogut gekannt haben wie \textsc{Ivory}. Nach dem art. 6. der \textit{Disquisitiones circa seriem} vom Jahre 1812, wo \textsc{Gauss} diese Entwicklung wiedergibt, hat man insbesondere für den von \(\theta\) freien Teil
\begin{equation}
\begin{aligned}
P&=\frac1{2\pi}\int_0^{2\pi}(a^2+b^2-2ab\cos\theta)^n\,d\theta\\
&=\frac1\pi\int_0^\pi(a^2+b^2-2ab\cos2\varphi)^n\,d\varphi
\end{aligned}\tag*{(4)}
\end{equation}
\par\bigskip\noindent\hfill67*\hspace*{.14\linewidth}
\par\endgroup

\clearpage\SourceLeaf{561}\ArithmeticHead{532}{L}{TAGEBUCH MIT ERLÄUTERUNGEN.}
\begingroup\small
\makeatletter\tagsleft@true\let\veqno\@@leqno\makeatother
\noindent die zwiefache Darstellung (siehe Werke III, S. 128, 129)
\begin{equation}
\begin{aligned}
P&=a^{2n}F\left(-n,-n,1,\frac{b^2}{a^2}\right)\\
&=(a+b)^{2n}F\left(-n,\frac12,1,\frac{4ab}{(a+b)^2}\right).
\end{aligned}\tag*{(5)}
\end{equation}
\noindent Mit freier Benutzung der \textsc{Ivory}schen Abhandlung läßt sich nun der Beweis für die Entwicklungen [1] und [2] wie folgt erbringen.

Bedeutet \(\varphi\) die exzentrische Anomalie, \(x\) die Exzentrizität einer Ellipse, deren große Achse gleich 1 ist, so hat man
\[
1-x^2\cos^2\varphi=1-\frac{x^2}{2}-\frac{x^2}{2}\cos2\varphi=a^2+b^2-2ab\cos2\varphi,
\]
\noindent für
\[
a=\frac{1+\sqrt{1-x^2}}2,\quad b=\frac{1-\sqrt{1-x^2}}2,\quad x^2=\frac{4ab}{(a+b)^2},
\]
\noindent so daß also \(x\) als das Verhältnis zwischen dem geometrischen und dem arithmetischen Mittel von \(a,b\) erscheint (\textsc{Landen}sche Transformation). Für \(n=\frac12\) ergibt sich aus (4) mit Benutzung der ersten Reihendarstellung in (5) für das vollständige elliptische Integral zweiter Gattung, d. h. für den halben Ellipsenumfang:
\[
\begin{aligned}
&\int_0^\pi(1-x^2\cos^2\varphi)^{\frac12}\,d\varphi\\
&=\int_0^\pi(a^2+b^2-2ab\cos2\varphi)^{\frac12}\,d\varphi\\
&=a\pi\biggl(1+\left(\frac12\right)^2\frac{b^2}{a^2}+\left(\frac{1\,.\,1}{2\,.\,4}\right)^2\frac{b^4}{a^4}\\
&\hspace{3em}+\left(\frac{1\,.\,1\,.\,3}{2\,.\,4\,.\,6}\right)^2\frac{b^6}{a^6}+\cdots\biggr).
\end{aligned}
\]
\noindent Setzt man hierin \(x=1\), also \(a=\frac12,\ \frac ba=1\), so erhält man die Reihe [1]. Für \(n=-\frac12\) und mit Benutzung der zweiten Reihendarstellung in (5) ergibt sich die von \textsc{Gauss} oft angewandte Entwicklung des vollständigen Integrals erster Gattung
\[
\begin{aligned}
&\int_0^\pi\frac{d\varphi}{\sqrt{1-x^2\cos^2\varphi}}\\
&=\int_0^\pi\frac{d\varphi}{\sqrt{a^2+b^2-2ab\cos2\varphi}}\\
&=\pi\biggl(1+\left(\frac12\right)^2x^2+\left(\frac{1\,.\,3}{2\,.\,4}\right)^2x^4\\
&\hspace{3em}+\left(\frac{1\,.\,3\,.\,5}{2\,.\,4\,.\,6}\right)^2x^6+\cdots\biggr),
\end{aligned}
\]
\noindent aus der die Gleichung
\[
\begin{aligned}
&\int_0^1d(x^2)\frac1\pi\int_0^\pi\frac{d\varphi}{\sqrt{1-x^2\cos^2\varphi}}\\
&\qquad=1+\left(\frac12\right)^2\cdot\frac12+\left(\frac{1\,.\,3}{2\,.\,4}\right)^2\cdot\frac13+\left(\frac{1\,.\,3\,.\,5}{2\,.\,4\,.\,6}\right)^2\cdot\frac14+\cdots
\end{aligned}
\]
\noindent hervorgeht. Linker Hand findet man nach Vertauschung der Integrationsfolge sofort den Wert \(\frac4\pi\), womit die Entwicklung [2] gewonnen ist. Die Annahme, daß das Verfahren von \textsc{Gauss} dem hier befolgten ähnlich gewesen sein mag, wird auch dadurch gestützt, daß im art. [3.] des Abschnitts [VII.] der \textit{Reihenlehre}, der von der Konvergenz der Entwicklung der Mittelpunktsgleichung handelt (siehe oben S. 423), die Entwicklung (3) für \(n=-\frac12\) benutzt wird.
\par\noindent\hfill\textsc{Schlesinger.}
\par\endgroup
